\documentclass[11pt]{article}

\usepackage[T1]{fontenc}
\usepackage[utf8]{inputenc}
\usepackage[english]{babel}
\usepackage[margin=2.3cm]{geometry}
\usepackage{amsmath,amssymb}
\usepackage{booktabs}
\usepackage{enumitem}
\usepackage{xcolor}
\usepackage{hyperref}
\usepackage{microtype}

\hypersetup{
  colorlinks=true,
  linkcolor=blue!60!black,
  urlcolor=magenta!70!black
}

\newcommand{\classification}[1]{\par\noindent\textbf{Classification: #1.}\par}
\newcommand{\location}[1]{\textbf{Location: #1}}
\newcommand{\problem}{\noindent\textbf{Problem.}\hspace{0.35em}}
\newcommand{\importance}{\par\noindent\textbf{Why it matters.}\hspace{0.35em}}
\newcommand{\fix}{\par\noindent\textbf{Suggested fix.}\hspace{0.35em}}

\title{Technical Review of\\
\emph{A High-Resolution Finite-Difference Benchmark for Steady Continuation and Hopf Bifurcation Dynamics up to Extreme Reynolds Numbers}}
\author{Manuscript review}
\date{23 September 2026}

\begin{document}
\maketitle

\section*{Overall assessment}

The manuscript has a useful objective and contains substantial computational work, but the present evidence does not support its central claims of a reference-grade high-Reynolds-number benchmark or of resolved Hopf, torus, and limit-cycle dynamics.  Major revision is required.  The highest-risk problems are not matters of presentation: the resumed unsteady histories are discontinuous because the archived flow state and archived telemetry correspond to different physical times; the reported Fourier frequencies are raw record-length bin frequencies; the hybrid-scheme definition and dissipation calculation are internally inconsistent; the GCI results are not reproducible from the reported values; and several validation and asymptotic claims contradict the manuscript's own tables or established reference values.

The governing nondimensional equations, the sign in the DST inversion, and the signs of the stated Thom wall formulas were checked and appear internally consistent.  The audit also checked branch conventions, sign conventions, symbol definitions, dimensional consistency, and implementation-facing expressions.  No inverse-trigonometric formula occurs in the manuscript.  The principal branch/sign failure is instead the signed-versus-unsigned cell-P\'eclet definition discussed below.

\section*{Major technical findings}

\subsection*{1. The resumed unsteady records contain a definite state--time mismatch}

\location{Unsteady-results section, Figs.~5--9; linked implementation in \path{unsteady_cavity_solver.py}.  Relevant routines are \path{capture_cycle_snapshots} and \path{resume_from_unsteady_npz}.}

\classification{Definite error}

\problem The telemetry plots show abrupt simultaneous jumps in probe velocity, kinetic energy, and/or enstrophy at concatenation boundaries (for example, near $t=10$ at $Re=15{,}000$, near $t=6$ and $12$ at $Re=25{,}000$, near $t=5$ and $10$ at $Re=50{,}000$, and near $t=3$ and $8$ at $Re=100{,}000$).  The linked implementation explains these jumps.  After \texttt{run\_simulation} stops at $t_{\mathrm{end}}$, \texttt{capture\_cycle\_snapshots} advances the solver through one additional inferred period.  The code then saves this advanced field as \texttt{final\_omega}/\texttt{final\_psi}, while the stored telemetry still ends at $t_{\mathrm{end}}$.  On resume, that advanced state is restored but assigned the last telemetry time.  Thus the continued record is not a continuous physical trajectory.

\importance The discontinuities contaminate the phase portraits, spectra, cycle counts, energy histories, snapshot phases, and all attractor classifications obtained from the concatenated signals.  They can themselves create low-frequency content and folded phase-space paths.

\fix Save the restart state immediately at the end of \texttt{run\_simulation}, before any diagnostic marching.  Generate phase snapshots from a deep copy of the solver, or advance and record the actual state time consistently.  Add a restart-continuity test requiring the first resumed sample to match a one-step uninterrupted reference run to numerical tolerance.  Regenerate every resumed unsteady dataset, figure, frequency, and conclusion after this correction.

\subsection*{2. The plotted records do not demonstrate stationary limit cycles or a 2-torus}

\location{Figs.~5--9 and the associated text on pp.~10--17.}

\classification{Unsupported claim}

\problem Even apart from the restart jumps, the global observables do not establish stationarity.  The kinetic energy in Fig.~5 continues to drift upward over the claimed stationary interval; Figs.~7--9 show substantial secular drift or step changes; and the marked entry and terminal points in the phase portraits generally do not coincide.  The $Re=15{,}000$ phase portrait is a non-closing multi-loop projection, but a two-dimensional projection with multiple loops is not sufficient evidence of an invariant 2-torus.  The abstract and conclusion additionally quote a recurrence error $\epsilon_{\mathrm{rec}}<0.08\%$, but $\epsilon_{\mathrm{rec}}$ is neither defined nor reported in the methods or results.

\importance A bounded-looking projection can arise from a transient, a restart artifact, slow drift, amplitude modulation, or quasiperiodicity.  The manuscript presently treats these alternatives as already excluded.

\fix After correcting the restart path, run uninterrupted histories until global observables and cycle-to-cycle diagnostics are stationary.  Define a phase-aligned recurrence norm, report it over several consecutive cycles, and show time-step and grid convergence.  For a torus claim, report two resolved incommensurate frequencies, a Poincar\'e section or suitable delay embedding, and stability under longer-window subdivision.  Otherwise use descriptive language such as ``nonstationary multi-frequency response'' rather than assigning a bifurcation topology.

\subsection*{3. The reported FFT frequencies are locked to record-length bins}

\location{Unsteady subsections and Table~5.}

\classification{Definite methodological error in the stated uncertainty claim}

\problem The reported frequencies coincide with integer Fourier bins of the declared analysis windows:

\begin{center}
\small
\begin{tabular}{@{}rrrrr@{}}
\toprule
$Re$ & Window length & $\Delta f=1/T_{\mathrm{win}}$ & Reported $St$ & Corresponding bin \\
\midrule
$10{,}000$  & $12.00$ & $0.08333$ & $0.6664$ & $8/12=0.6667$ \\
$25{,}000$  & $17.80$ & $0.05618$ & $0.1124$ & $2/17.8=0.11236$ \\
$50{,}000$  & $16.67$ & $0.05999$ & $0.2399$ & $4/16.67=0.23995$ \\
$100{,}000$ & $13.20$ & $0.07576$ & $0.2272$ & $3/13.2=0.22727$ \\
\bottomrule
\end{tabular}
\end{center}

In particular, observing only two or three cycles does not ``permanently eliminate'' sampling-window or bin artifacts; it makes frequency uncertainty large.  A window function reduces leakage but does not create the missing frequency resolution.

\importance The quoted four-decimal frequencies and derived periods imply precision far beyond that supported by the records.  The apparent high-$Re$ clustering near $St\approx0.23$ can be a consequence of selecting nearby low-order bins.

\fix Use substantially longer uninterrupted stationary records, state the spectral estimator and detrending/windowing procedure, and report uncertainty.  Demonstrate consistency across disjoint time windows and with a time-domain period estimate.  Until then, report at most the occupied FFT bin and its resolution, not four-decimal physical frequencies.

\subsection*{4. The $Re=15{,}000$ case is internally contradictory}

\location{The ``Secondary Hopf Bifurcation'' subsection, Fig.~6, abstract, conclusion, and Table~5.}

\classification{Definite error}

\problem The subsection states $25{,}000$ steps with $\Delta t=4\times10^{-4}$, hence $t=10$, and reports $St=0.6662$ and $T=1.5010$.  Figure~6 extends to $t=20$ and labels its spectral peak $St=0.5554$, $T=1.8004$; Table~5, the abstract, and the conclusion also use $t_{\mathrm{end}}=20$, $St=0.5554$, and $T=1.8004$.  The figure itself is titled ``Hopf Limit-Cycle Attractor'', while the caption and prose call the same object a 2-torus.

\importance These are mutually incompatible descriptions of the run and its principal result.

\fix Identify the authoritative dataset, regenerate the subsection and all summary locations from it, and state the exact uninterrupted analysis window.  Do not classify the attractor until the evidence requested in Finding~2 is supplied.

\subsection*{5. The hybrid discretization, P\'eclet notation, and numerical-viscosity audit do not describe one consistent scheme}

\location{Cell-P\'eclet and localized-hybrid subsection, Eqs.~(18)--(20), Fig.~1, abstract, and conclusion.}

\classification{Definite formulation and arithmetic errors}

\problem First, $Pe_{x,i,j}$ is defined as $|u_{i,j}|h/\nu\ge0$, but the piecewise derivative later contains the branch $Pe_{x,i,j}<-2$.  That branch is unreachable under the stated definition.  A signed quantity such as $\widetilde{Pe}_x=u h/\nu$ is required for directional branching, with $|\widetilde{Pe}_x|$ used only for the threshold.

Second, the printed dissipation equation does not yield the printed values.  With the manuscript's core values $u=-0.0304$, $v=-0.0036$, $h=1/512$, and $\nu=10^{-5}$,
\[
\frac{1}{2\nu}\left[(|u|h-2\nu)_+ + (|v|h-2\nu)_+\right]=1.97,
\]
not $0.99$.  At $|u|=1$ it gives $96.66$, not $48.33$.  The plotted/reported values therefore contain an additional factor of two that is absent from the equation.

Third, the modified-equation numerical diffusion of the explicitly written first-order upwind derivative is $|u|h/2$ in the $x$ direction (and $|v|h/2$ in the $y$ direction).  That is not the same as subtracting $\nu$ after the $Pe=2$ switch.  The manuscript appears to mix a clipped-coefficient hybrid interpretation with a piecewise replacement by a full upwind derivative.

\importance The $Re=100{,}000$ accuracy and dissipation claims depend directly on this audit.  Moreover, the method is described as ``localized'', although the manuscript says upwinding is active on $98.21\%$ of the domain.  That is nearly domain-wide activation, not localization to corner shocks or wall layers.

\fix State the exact discrete operator implemented in code using a signed P\'eclet number.  Derive its leading modified-equation diffusion component by component, define whether a scalar plotted value is a sum, average, maximum, or tensor diagnostic, and recompute Fig.~1.  Replace ``localized'' with an accurate description unless the switching rule is actually spatially restricted.  Reassess every $Re=100{,}000$ benchmark after correcting this definition.

\subsection*{6. The GCI claims are not reproducible and do not establish asymptotic convergence}

\location{Formal GCI subsection, Table~3, and Fig.~3.}

\classification{Definite reproducibility failure; unsupported conclusion}

\problem Applying the manuscript's own equations to the displayed $Re=10{,}000$ values gives approximately
\[
p=4.7285,\qquad f_{h=0}=-0.1286984,\qquad \mathrm{GCI}_{21}=0.0179\%,
\]
not the tabulated $-0.12871$ and $0.031\%$.  Unrounded data may change these results, but those data and the observed order are not reported.  At $Re=100{,}000$, $f_1=f_2=-0.12572$ to all displayed digits, so $p$ and the quoted ${<}0.001\%$ GCI cannot be reconstructed at all.  The very high apparent order at $Re=10{,}000$ also requires an asymptotic-range or cancellation check, especially because the wall closure is first order and the $Re=100{,}000$ operator switches scheme with grid spacing.  Figure~3 shows only two grids, not the complete triplet used by the calculation.

\importance A small change in one scalar extremum is not sufficient to establish mesh independence of the full field, wall layers, vortex topology, or unsteady dynamics.

\fix Report full-precision $f_i$, the observed order $p$, convergence ratio, iterative residuals, and both coarse/fine GCIs.  Demonstrate the asymptotic-range consistency check and include all three grids.  Apply convergence studies to centerline norms, vorticity, wall quantities, vortex centers/strengths, and unsteady frequencies, not only $\psi_{\min}$.

\subsection*{7. The validation prose overstates the agreement shown by the tables and figures}

\location{Table~1 and its discussion; Table~4, Fig.~4, and the near-lid validation discussion; abstract and conclusion.}

\classification{Unsupported claim}

\problem Table~1 gives $\psi_{\min}=-0.12868$ versus $-0.11970$ at $Re=10{,}000$ (about $7.5\%$ difference in magnitude), and $-0.12631$ versus $-0.11780$ at $Re=25{,}000$ (about $7.2\%$).  These differences are omitted when the text calls the agreement outstanding or exceptional.  The near-lid table is more problematic: at $Re=25{,}000$, the mid-plane values $-0.0085$ and $-1.9814$ differ by about $99.6\%$ relative to the reference magnitude; at $Re=30{,}000$ and $x=0.1563$, $5.4027$ versus $2.1243$ is a $154\%$ relative discrepancy.  The reported MAE exceeds $1$ at four of seven Reynolds numbers, and Fig.~4 visibly shows large high-$Re$ deviations and endpoint spikes.  These results do not constitute ``extraordinary fidelity'' or verification of zero numerical dissipation.

\importance Selective emphasis on one favorable station obscures the overall validation quality and weakens trust in the benchmark designation.

\fix Report normalized error norms, maximum error and its location, and uncertainty across all stations.  Discuss the significant discrepancies directly, avoid converting one local agreement into a global percentage claim, and replace promotional adjectives with quantitative statements.  Resolve whether the compared solutions use the same steady branch, boundary treatment, coordinates, and interpolation protocol.

\subsection*{8. The claimed Batchelor plateau value is likely inconsistent with the cited square-cavity limit}

\location{Asymptotic-core subsection, abstract, and conclusion.}

\classification{Likely issue requiring external verification}

\problem Batchelor's theorem supports uniform vorticity in a closed-streamline core; it does not by itself validate the manuscript's value $\omega_0=-2.5418$.  Under the manuscript's stated nondimensionalization and vorticity definition, standard square-cavity results approach approximately $-1.886$ (sign adjusted for the present clockwise convention), with benchmark values near $-1.9$ by $Re=10{,}000$.  The manuscript's values near $-2.54$ are therefore not a small deviation.  A standard deviation of $0.10$ over one centerline segment also does not demonstrate a two-dimensional plateau.

\importance The manuscript presents this number as confirmation of an asymptotic theorem, but it may instead indicate a vorticity scaling, differentiation, state-selection, or post-processing error.

\fix Verify $\omega$ independently from both $-\nabla^2\psi$ and $\partial v/\partial x-\partial u/\partial y$, document the normalization, compare the vortex-center and geometric-center values, and show a two-dimensional core mask with mean, spread, and grid dependence.  Reconcile the result explicitly with Batchelor/Burggraf and with the values in the cited cavity benchmarks before retaining the theorem-confirmation claim.

\subsection*{9. Physical-time accuracy and the high-$Re$ steady branch need clearer verification}

\location{ADI subsection, wall-vorticity subsection, homotopy subsection, and all unsteady results.}

\classification{Likely issue and implementation ambiguity}

\problem The nonlinear convection term is evaluated explicitly at level $n$ in both half-steps, so the temporal order of the complete method is not established by the diffusion-only Peaceman--Rachford structure.  No time-step convergence study is given.  At $Re=100{,}000$, the physical run additionally uses wall-vorticity relaxation $\beta_{\mathrm{wall}}=0.75$, which is a temporal boundary filter rather than the stated instantaneous no-slip closure.  The linked implementation also auto-selects relaxed wall updates at high $Re$.  Separately, the paper reports ``steady'' solutions above the onset of physical instability without giving a governing-equation residual, stopping criterion, or explanation of how an unstable stationary branch is isolated by pseudo-time marching.

\importance Frequencies, bifurcations, and steady-branch benchmarks can be dominated by temporal truncation, relaxation, or numerical stabilization.  A $513^2$ grid with only about $1.62$ cells across the nominal $Re=100{,}000$ boundary-layer thickness cannot be called resolved without a targeted convergence demonstration.

\fix State the formal temporal order of the implemented scheme and verify frequencies and amplitudes under at least two smaller $\Delta t$ values.  For physical runs, use $\beta_{\mathrm{wall}}=1$ or quantify and extrapolate the relaxation effect.  For stationary branches, report maximum discrete residuals, iteration histories, and the continuation/stabilization procedure, and call them stationary numerical branches rather than physically stable steady flows where appropriate.

\section*{Notation and definition audit}

\begin{itemize}[leftmargin=*]
  \item $Pe_x$ and $Pe_y$ are first defined as nonnegative magnitudes but are later used as signed directional quantities.  This is a direct descriptor and implementation conflict; introduce separate signed and magnitude notation.
  \item $St$ at $Re=15{,}000$ changes from the subsection's ``carrier frequency'' $0.6662$ to the table/figure's ``dominant frequency'' $0.5554$ without distinguishing carrier, modulation, or sideband frequencies.
  \item $\epsilon_{\mathrm{rec}}$ appears in the abstract and conclusion without a definition, norm, phase alignment, or data table.
  \item The manuscript alternates between ``geometric center'', ``primary vortex center'', and ``core'' when assigning $\omega_0$.  These are not interchangeable measurement locations unless the two-dimensional plateau has first been demonstrated.
  \item ``Steady'', ``converged'', and ``physical time-accurate'' are used without corresponding residual or temporal-convergence definitions.  Define these terms operationally.
\end{itemize}

\section*{Meaningful editorial and figure-consistency issues}

\begin{enumerate}[leftmargin=*]
  \item \location{Abstract and Introduction.} The statement that established benchmarks ``stagnate at $Re\le10{,}000$'' conflicts with the manuscript's own use of Erturk et al. at higher Reynolds numbers.  Replace it with a precise statement about which fields, grids, or formats are missing.
  \item \location{Unsteady figures.} The generated figure titles label the $Re=15{,}000$ through $100{,}000$ cases as ``Hopf Limit Cycle Onset'' or ``Limit-Cycle Attractor'' even where the manuscript calls them a torus, folded dynamics, or a generic bounded attractor.  Regenerate titles from the final classifications rather than hard-coding one label.
  \item \location{Near-lid validation.} ``Stepping just two grid cells away'' is inconsistent with the coordinates quoted: on a $513^2$ grid, $x=0.9688$ is roughly sixteen cells from $x=1$, while $x=0.9800$ is not a native grid coordinate.  State whether these are benchmark stations or interpolated points.
  \item \location{Numerical-method and data-availability prose.} Markdown backticks are used around \texttt{scipy.fft.dst} and \texttt{.npz}; use \verb|\texttt{...}| in LaTeX.  Likewise, format \texttt{kx=3, ky=3} consistently as code or as mathematical subscripts.
\end{enumerate}

\section*{Proofing sweep}

The bundled mechanical scan was run on the compiled manuscript.  Its only pattern hit was the legitimate mathematical ellipsis in ``$1,\ldots,N-2$'', so no duplicated-punctuation defect was confirmed.  A manual spot-check of the bibliography found no duplicated commas, malformed quotation punctuation, or broken citation keys.  The following small corrections remain worthwhile:

\begin{itemize}[leftmargin=*]
  \item p.~4: replace ``analytical sinus-eigenmodes'' with ``analytical sine eigenmodes'' or ``sinusoidal eigenmodes'';
  \item throughout: use typographic ranges, for example $Re=10{,}000$--$30{,}000$, rather than a mathematical minus sign;
  \item keep code identifiers in \verb|\texttt{}| and reserve math italics for mathematical variables.
\end{itemize}

\section*{Items requiring external verification}

\begin{enumerate}[leftmargin=*]
  \item The cited Erturk--Corke--G\"ok\c{c}\"ol study describes steady solutions for $Re<21{,}000$, yet Table~1 attributes a $Re=25{,}000$ reference value to \texttt{erturk2005numerical}.  Verify the actual source of that row and correct the citation.
  \item The cited Peng--Shiau--Hwang study reports a first Hopf threshold near $Re=7402\pm4\%$ and a frequency near $0.59$.  Reconcile the manuscript's $Re_{\mathrm{cr}}\approx8{,}000$--$10{,}000$ and $St=0.6664$ with the cited definitions, lid treatment, and numerical method rather than saying the cited work documents the same result.
  \item The linked GitHub README, accessed on 23 September 2026, lists unsteady frequencies that differ substantially from the manuscript (for example $St=0.2856$, $0.3331$, and $0.5128$ at $Re=25{,}000$, $50{,}000$, and $100{,}000$, respectively).  Version the code and data used for the paper, record a commit hash, and make the manuscript, archive, README, and generated figures agree.
  \item Confirm that the Zenodo record contains the claimed twenty-eight datasets, the exact scripts used to produce every table and figure, environment metadata, and checksums tied to the reviewed manuscript version.
\end{enumerate}

\section*{Highest-priority fixes before resubmission}

\begin{enumerate}[leftmargin=*]
  \item Correct the state--time restart bug and regenerate all unsteady datasets and figures from continuous histories.
  \item Re-establish stationarity, temporal/grid convergence, frequency uncertainty, and attractor classification; otherwise remove the Hopf/torus/limit-cycle claims.
  \item Rewrite and rederive the hybrid operator and numerical-dissipation analysis, including Fig.~1 and every $Re=100{,}000$ conclusion.
  \item Recompute the GCI study from full-precision values and report residuals and observed orders.
  \item Recalibrate the benchmark and Batchelor claims to the actual errors and accepted reference values.
  \item Synchronize the abstract, body, tables, figures, conclusion, repository, and archive from one versioned set of results.
\end{enumerate}

Until these points are resolved, the manuscript should not describe the data as a validated reference benchmark or the unsteady trajectories as established bifurcation attractors.

\end{document}